% Chapter 2, section 8: Headers and routing
% Write in plain LaTeX. Code of a guide step: \guide{<step>}; commands: \begin{shell}...\end{shell};
% a screenshot: \screenshot{file}{caption}; a gap in the draft: \placeholder{...}.
% See book/latex/queen.sty for the full list.
%
% Draft by Claude, 2026-11-16, from the website's headers and routing pages
% (webdoc/src/content/docs/writing-applications/headers.mdx, routing.mdx). The code is cut from
% examples/tutorials/js/book-client-guide.mjs, which ran it against a 2.9.0-beta.3 node and checked
% every result quoted here, and the refusals in the tables marked by their code (where_mismatch,
% reading_queue, has_readers, duplicate, source_mismatch).

\section{Headers and routing}

A message can carry headers beside its payload. A consumer group can take only some of a queue's
messages, a queue can read other queues, and a transaction can forward the message it holds into
another queue. In each case the message is written once, at its source: the group or the queue
that reads it holds a \emph{link}, about 20 bytes, written in the same commit as the message.

Three rules are all there is to learn:

\begin{enumerate}
\item A group or a queue reads only what its \texttt{where} takes.
\item The source keeps the messages; the queues that read it consume them.
\item It is the same message, with the same \texttt{id} and \texttt{transactionId}, everywhere it
shows.
\end{enumerate}

\why{without this, a reader that wants a few kinds of message reads and acks all of them, and
routing means a router process that copies every message and adds a commit for each batch. A link
costs the source's commit a few bytes, and it cannot be lost or doubled: it is in the same log
entry as the message.}

The examples read one queue, \texttt{shop-events}, where a shop writes the events of its orders,
one partition per order. It is configured like any other queue.

All of this is new in QueenMQ 2.9. Every node of the cluster must run it: before that, headers, a
\texttt{where}, \texttt{from} or a link are refused with \texttt{cluster\_version}.

\subsection{Headers}

Headers are a small map of strings stored beside the payload: what the message is, where it comes
from, which trace it belongs to.

\guide{headers}

They come back on every read of the message: a pop, a fetch, the feed, the message browser and the
dead-letter queue.

\why{without headers, the type and the trace context go inside the payload. Then every payload must
be a JSON object, every filter parses the body, and a tool that must not read the body (it may hold
personal data) has nothing to read.}

\begin{pairs}[0.20]{What}{Rule}
Names & 1 to 64 bytes of \texttt{a-z}, \texttt{0-9}, \texttt{.}, \texttt{\_} and \texttt{-}: lower
case only \\
Values & any UTF-8 string \\
Reserved & names that start with \texttt{queen-} are the broker's \\
How many & 64 per message (\texttt{QUEEN\_MAX\_HEADERS}) \\
How big & 8 KiB per message, names and values together (\texttt{QUEEN\_MAX\_HEADER\_BYTES}, at most
64 KiB) \\
Encryption & never: headers are stored in clear on an encrypted queue too, so routing and
monitoring read them without the body. Put no secret and no personal data in a header. \\
Dedup & not part of it: the \texttt{transactionId} stays the key \\
Dead letters & a dead letter keeps its headers, and a replay sends them again \\
Timers & a timer takes \texttt{headers} too, and the message it fires carries them \\
Kafka & a record produced through the Kafka port keeps its headers; the ones whose value is UTF-8
are also the message's headers, with their names in lower case \\
\end{pairs}

A header describes the message. It never says where the message goes: QueenMQ routes on a header
only when a reader asks for it, in its own \texttt{where}. Some names that tools understand, documented
but not enforced: \texttt{traceparent} and \texttt{tracestate} (W3C trace context), \texttt{type},
\texttt{content-type}.

\subsection{The where language}

Filtered groups, reading queues and the tenant feed filter with one language. A \texttt{where} can be
one string:

\guide{where-text}

or a list of conditions, which is what \texttt{.where(field, op, value)} builds.

\begin{pairs}[0.20]{Part}{What it takes}
Fields & \texttt{queue}, \texttt{partition}, \texttt{transactionId}, \texttt{data} or
\texttt{data.a.b} (a path into the JSON payload, up to 16 deep), \texttt{headers.<name>} \\
Operators & \texttt{=} \texttt{!=} \texttt{<} \texttt{<=} \texttt{>} \texttt{>=} \texttt{in}
\texttt{exists} \texttt{prefix} \texttt{!prefix} \\
Combining & \texttt{and} only; \texttt{in (...)} for alternatives \\
Values & \texttt{'text'}, numbers, \texttt{true}, \texttt{false}, \texttt{null}; \texttt{in} takes up
to 64 values \\
Size & up to 16 conditions \\
\end{pairs}

A field the message does not have holds for no operator, \texttt{!=} included, as in SQL. The
\texttt{where} above took the Italian and French events over 100, and not the German order.

The leader evaluates a \texttt{where} once, when it plans the push, and the log entry keeps the
answer: followers and replays never evaluate it again. On an encrypted queue, the leader decrypts
the payload for it.

\subsection{Filtered consumer groups}

A consumer group with a \texttt{where} gets only the messages it takes. Billing needs the payments
and the refunds, not every event of an order:

\guide{filtered-group}

The group got four events of the six, with their headers, in the order of their partitions. Its
\texttt{where} is set at its first pop, like its start, and stays fixed from then on. The group
gets only its messages, its lag counts only those, and a waiting pop wakes only for one of them.
Leases, acks, retries and dead letters work as for any group, with the queue's settings.

\begin{pairs}[0.40]{Case}{Answer}
Another worker of the group sends another \texttt{where} & refused, \texttt{where\_mismatch}, with
the group's current \texttt{where} \\
A \texttt{where} without a group & refused: name a group \\
A \texttt{where} on a group that already reads the queue without one & refused,
\texttt{where\_mismatch}: a filter starts as a new group \\
\texttt{subscriptionMode('all')} or \texttt{subscriptionFrom} & the group starts with every match
the queue still keeps \\
Change the \texttt{where} & the admin call \texttt{updateConsumerGroupWhere(group, queue, where)}:
it applies from the next push, and a worker that still sends the old one is refused \\
Delete the group & as any group, with the admin call \texttt{deleteConsumerGroupForQueue(group, queue)} \\
The group listings & show the group under its queue, with \texttt{filtered: true} and its
\texttt{where} \\
\end{pairs}

\why{``only some of these'' is what most readers need. Here it is one line in the consumer, with no
queue to create, no router to run and nothing to ack that the consumer does not want.}

\subsection{Reading queues}

A reading queue is a queue whose messages come from other queues. Give it \texttt{from} and, if you
want, \texttt{where}, \texttt{partitionBy} and \texttt{start}:

\guide{reading-queue}

The loyalty service gets the payments and the refunds too, but by customer: partition
\texttt{customer-42} holds the events of every order of that customer, in the order the log
committed them, each with the \texttt{id} it has at the source.

\begin{options}{Option}{Default}{What it does}
from & (none) & The source: a queue name, or a list of names and \texttt{\{ queue, where \}}
objects. \texttt{null} stops the reading; what the queue holds stays readable. \\
where & every message & The condition a message of a source named by its name alone must meet. \\
partitionBy & \texttt{same} & \texttt{same}: the source's partitions, in the source's order.
\texttt{data.<path>} or \texttt{headers.<name>}: the field's value names the partition (a message
without it is not routed). \\
start & \texttt{now} & \texttt{now}: the messages pushed from now on. \texttt{all}: every match the
sources still keep, then the new ones, in order. \\
\end{options}

A reading queue has its own leases, retries, dead-letter queue, priority, delays and consumer
groups, and a token of its own can read it without reading its sources. It has no keep settings:
retention, the dedup window, encryption, \texttt{maxSize} and the sink hold are refused on it
(\texttt{reading\_queue}), because its sources decide how long a message is kept. A push into it is
refused too: push to its source.

\why{a filtered group shares its queue's settings and its partitions. A reading queue is for a
reader that needs settings of its own, partitions of its own (\texttt{partitionBy}), several
sources, or a token of its own.}

\begin{pairs}[0.40]{Case}{Answer}
An ordinary queue given \texttt{from} & it becomes a reading queue if no message was ever pushed to
it; refused once it holds messages \\
A reading queue that reads itself, at any depth & refused, \texttt{route\_loop} \\
A source that does not exist & refused, \texttt{source\_not\_found} \\
Two sources route one message (one \texttt{transactionId}) into one partition & kept once: a reading
queue deduplicates by \texttt{transactionId} in its partitions, as every queue does \\
Change its \texttt{where} & from the next push; to rebuild it, delete it and create it again \\
Delete a source that reading queues read & refused, \texttt{has\_readers} (409), unless
\mbox{\texttt{delete(\{ force: true \})}}: their links then read as gone \\
\end{pairs}

\subsection{Forward a message: link}

A transaction can put the message it holds into another queue without sending its payload again,
together with the ack, the KV writes and the timers, all or nothing. Here a fraud check sends the
risky payments to a person:

\guide{link}

The target gets the same message, with the same \texttt{id} and \texttt{transactionId}, and the
source's headers with the link's own written over them. Without \texttt{.partition(p)}, the message
keeps the name of its partition. The same link again commits and writes nothing (its result is
marked \texttt{duplicate}), so a forward is safe to retry. A link to a message that is gone rolls the
transaction back (\texttt{source\_gone}), and so does one whose \texttt{transactionId} is not the one
at that offset (\texttt{source\_mismatch}).

\why{a pipeline that hands a message to the next step sends the whole payload again at every hop.
A link costs about 20 bytes, whatever the payload weighs, and it commits with the ack of the input,
so the hand-over cannot be lost or done twice.}

A stream's \texttt{.to(queue)} writes a link too, when its result is its input unchanged (nothing
but \texttt{where}, \texttt{filter}, \texttt{keyBy} or routing before it); with the option
\mbox{\texttt{link: false}} it writes copies.

\subsection{What the readers keep, and their lag}

The source keeps the message. For the source's retention, every filtered group, every reading
queue and every queue that holds its links counts as one more consumer group:

\begin{pairs}[0.40]{Source setting}{What a reader gets}
\texttt{retentionSeconds}, \texttt{maxWaitTimeSeconds} & the same limit: past it, a link reads as
gone, as a message is gone for a group that is behind \\
\texttt{completedRetentionSeconds} & the source waits for the readers too: it keeps a message until
they have consumed their link \\
retention off & everything is kept \\
\end{pairs}

So a queue that holds links keeps the source's messages for as long as it keeps the links: give a
queue you forward into a retention of its own. The source's page in the dashboard, and
\texttt{GET /api/v1/resources/queues/:queue}, list its readers and its holders. A dead letter filed
from a reading queue is a copy: it holds the message itself, because a dead letter outlives its
source.

A filtered group's lag counts only the messages its \texttt{where} took. A message that no reader
takes waits for nobody: it counts as completed. The partitions no reader can take at all (when the
\texttt{where}s name queues or partitions) show as \emph{unread}, and the dashboard marks the queue
``No reader'' for them.

\why{a filtered group that counted the whole queue as its lag would look behind for ever, and an
autoscaler reading its depth would start workers for messages it never gets. The partitions nobody
reads are the gap worth seeing: a worker that was never started.}

\subsection{Limits and traps}

\begin{options}{Setting}{Default}{What it limits}
QUEEN\_WHERE\_MAX\_CONDITIONS & 16 & The conditions in one \texttt{where} \\
QUEEN\_WHERE\_MAX\_IN & 64 & The values in one \texttt{in} \\
QUEEN\_MAX\_READERS\_PER\_QUEUE & 64 & The readers of one source: reading queues and filtered groups \\
QUEEN\_MAX\_SOURCES\_PER\_QUEUE & 16 & The sources of one reading queue \\
\normalfont(fixed) & 8 deep & Readers of readers, in a chain \\
\end{options}

Every reader costs an evaluation of its \texttt{where} on each push to its source. That is why the
readers of one source are limited; raise the limit knowing it.

\begin{itemize}
\item \textbf{A new filtered group starts at the end of the queue}, like any group: ask for
\mbox{\texttt{subscriptionMode('all')}} to read what is already there.
\item \textbf{A \texttt{where} is fixed at the group's first pop.} To change it, use the admin
call in the table of filtered groups, then deploy the workers with the new one.
\item \textbf{Headers are in clear}, on an encrypted queue too.
\item \textbf{A queue you forward into holds the source's messages}: give it retention.
\end{itemize}

On the site: \href{https://queenmq.com/writing-applications/headers/}{queenmq.com/writing-applications/headers}
and \href{https://queenmq.com/writing-applications/routing/}{queenmq.com/writing-applications/routing}.
